%% Generated by Sphinx.
\def\sphinxdocclass{report}
\documentclass[letterpaper,10pt,english]{sphinxmanual}
\ifdefined\pdfpxdimen
   \let\sphinxpxdimen\pdfpxdimen\else\newdimen\sphinxpxdimen
\fi \sphinxpxdimen=.75bp\relax
\ifdefined\pdfimageresolution
    \pdfimageresolution= \numexpr \dimexpr1in\relax/\sphinxpxdimen\relax
\fi
%% let collapsible pdf bookmarks panel have high depth per default
\PassOptionsToPackage{bookmarksdepth=5}{hyperref}

\PassOptionsToPackage{booktabs}{sphinx}
\PassOptionsToPackage{colorrows}{sphinx}

\PassOptionsToPackage{warn}{textcomp}
\usepackage[utf8]{inputenc}
\ifdefined\DeclareUnicodeCharacter
% support both utf8 and utf8x syntaxes
  \ifdefined\DeclareUnicodeCharacterAsOptional
    \def\sphinxDUC#1{\DeclareUnicodeCharacter{"#1}}
  \else
    \let\sphinxDUC\DeclareUnicodeCharacter
  \fi
  \sphinxDUC{00A0}{\nobreakspace}
  \sphinxDUC{2500}{\sphinxunichar{2500}}
  \sphinxDUC{2502}{\sphinxunichar{2502}}
  \sphinxDUC{2514}{\sphinxunichar{2514}}
  \sphinxDUC{251C}{\sphinxunichar{251C}}
  \sphinxDUC{2572}{\textbackslash}
\fi
\usepackage{cmap}
\usepackage[T1]{fontenc}
\usepackage{amsmath,amssymb,amstext}
\usepackage{babel}



\usepackage{tgtermes}
\usepackage{tgheros}
\renewcommand{\ttdefault}{txtt}



\usepackage[Sonny]{fncychap}
\ChNameVar{\Large\normalfont\sffamily}
\ChTitleVar{\Large\normalfont\sffamily}
\usepackage{sphinx}

\fvset{fontsize=auto}
\usepackage{geometry}


% Include hyperref last.
\usepackage{hyperref}
% Fix anchor placement for figures with captions.
\usepackage{hypcap}% it must be loaded after hyperref.
% Set up styles of URL: it should be placed after hyperref.
\urlstyle{same}

\addto\captionsenglish{\renewcommand{\contentsname}{محتوای مستندات:}}

\usepackage{sphinxmessages}
\setcounter{tocdepth}{1}



\title{Kunqalay}
\date{آوریل 30, 2025}
\release{1.0}
\author{Documentation Team}
\newcommand{\sphinxlogo}{\vbox{}}
\renewcommand{\releasename}{انتشار}
\makeindex
\begin{document}

\ifdefined\shorthandoff
  \ifnum\catcode`\=\string=\active\shorthandoff{=}\fi
  \ifnum\catcode`\"=\active\shorthandoff{"}\fi
\fi

\pagestyle{empty}
\sphinxmaketitle
\pagestyle{plain}
\sphinxtableofcontents
\pagestyle{normal}
\phantomsection\label{\detokenize{index::doc}}



\chapter{معرفی}
\label{\detokenize{index:id2}}
\sphinxAtStartPar
به مستندات پروژه کانقالای خوش آمدید! این مستندات به منظور ارائه اطلاعات جامع درباره ساختار، نحوه کارکرد و استفاده از سیستم کانقالای تهیه شده است.

\sphinxAtStartPar
کانقالای یک سیستم مدیریت منابع سازمانی (ERP) و مدیریت ارتباط با مشتری (CRM) است که با استفاده از فریم‌ورک Django پیاده‌سازی شده و امکانات متنوعی از جمله مدیریت مشتریان، مدیریت فروش، مدیریت محصولات، ارسال پیامک و پیام‌های واتساپ را فراهم می‌کند.


\chapter{بخش‌های مستندات}
\label{\detokenize{index:id3}}
\sphinxstepscope


\section{معرفی}
\label{\detokenize{introduction/index:introduction}}\label{\detokenize{introduction/index:id1}}\label{\detokenize{introduction/index::doc}}
\sphinxAtStartPar
در این بخش به معرفی کلی پروژه کانقالای می‌پردازیم.


\subsection{پروژه کانقالای}
\label{\detokenize{introduction/index:id2}}
\sphinxAtStartPar
کانقالای یک سیستم مدیریت منابع سازمانی (ERP) و مدیریت ارتباط با مشتری (CRM) است که برای مدیریت عملیات یک کسب و کار طراحی شده است. این سیستم امکان مدیریت مشتریان، فروش، محصولات، و ارتباطات را فراهم می‌کند.


\subsection{قابلیت‌های کلیدی}
\label{\detokenize{introduction/index:id3}}\begin{itemize}
\item {} 
\sphinxAtStartPar
مدیریت مشتریان (CRM)

\item {} 
\sphinxAtStartPar
مدیریت منابع سازمانی (ERP)

\item {} 
\sphinxAtStartPar
سیستم پیامک و ارتباط با مشتری

\item {} 
\sphinxAtStartPar
سیستم پیام‌رسانی واتساپ

\item {} 
\sphinxAtStartPar
مدیریت فروش و سفارش‌ها

\item {} 
\sphinxAtStartPar
مدیریت محصولات

\item {} 
\sphinxAtStartPar
گزارش‌گیری و داشبورد مدیریتی

\end{itemize}


\subsection{فناوری‌های مورد استفاده}
\label{\detokenize{introduction/index:id4}}\begin{itemize}
\item {} 
\sphinxAtStartPar
Django: فریم‌ورک اصلی توسعه وب

\item {} 
\sphinxAtStartPar
Celery: سیستم اجرای وظایف غیرهمزمان

\item {} 
\sphinxAtStartPar
Redis: سیستم کش و پردازش پیام‌ها

\item {} 
\sphinxAtStartPar
SQLite: پایگاه داده

\item {} 
\sphinxAtStartPar
Django Channels: برای ارتباطات WebSocket و ASGI

\item {} 
\sphinxAtStartPar
REST Framework: برای توسعه API

\end{itemize}

\sphinxstepscope


\section{راهنمای نصب و راه‌اندازی}
\label{\detokenize{installation/index:installation}}\label{\detokenize{installation/index:id1}}\label{\detokenize{installation/index::doc}}
\sphinxAtStartPar
در این بخش، مراحل نصب و راه‌اندازی پروژه کانقالای توضیح داده شده است.


\subsection{پیش‌نیازها}
\label{\detokenize{installation/index:id2}}
\sphinxAtStartPar
برای نصب و راه‌اندازی پروژه کانقالای، نیاز به موارد زیر دارید:
\begin{itemize}
\item {} 
\sphinxAtStartPar
پایتون 3.8 یا بالاتر

\item {} 
\sphinxAtStartPar
Redis

\item {} 
\sphinxAtStartPar
SQLite (یا دیگر پایگاه داده‌های پشتیبانی شده توسط Django)

\item {} 
\sphinxAtStartPar
pip (برای نصب وابستگی‌ها)

\end{itemize}


\subsection{مراحل نصب}
\label{\detokenize{installation/index:id3}}\begin{enumerate}
\sphinxsetlistlabels{\arabic}{enumi}{enumii}{}{.}%
\item {} 
\sphinxAtStartPar
\sphinxstylestrong{کلون کردن مخزن پروژه}

\sphinxAtStartPar
ابتدا مخزن پروژه را از گیت‌هاب کلون کنید:

\begin{sphinxVerbatim}[commandchars=\\\{\}]
git\PYG{+w}{ }clone\PYG{+w}{ }https://github.com/your\PYGZhy{}username/kunqalay.git
\PYG{n+nb}{cd}\PYG{+w}{ }kunqalay
\end{sphinxVerbatim}

\item {} 
\sphinxAtStartPar
\sphinxstylestrong{ایجاد محیط مجازی پایتون}

\sphinxAtStartPar
برای جلوگیری از تداخل با دیگر پروژه‌ها، توصیه می‌شود یک محیط مجازی ایجاد کنید:

\begin{sphinxVerbatim}[commandchars=\\\{\}]
python\PYG{+w}{ }\PYGZhy{}m\PYG{+w}{ }venv\PYG{+w}{ }venv
\PYG{c+c1}{\PYGZsh{} در ویندوز}
venv\PYG{l+s+se}{\PYGZbs{}S}cripts\PYG{l+s+se}{\PYGZbs{}a}ctivate
\PYG{c+c1}{\PYGZsh{} در لینوکس/مک}
\PYG{n+nb}{source}\PYG{+w}{ }venv/bin/activate
\end{sphinxVerbatim}

\item {} 
\sphinxAtStartPar
\sphinxstylestrong{نصب وابستگی‌ها}

\sphinxAtStartPar
با استفاده از pip، وابستگی‌های پروژه را نصب کنید:

\begin{sphinxVerbatim}[commandchars=\\\{\}]
pip\PYG{+w}{ }install\PYG{+w}{ }\PYGZhy{}r\PYG{+w}{ }requirements.txt
\end{sphinxVerbatim}

\item {} 
\sphinxAtStartPar
\sphinxstylestrong{اجرای مهاجرت‌های پایگاه داده}

\sphinxAtStartPar
پایگاه داده را با اجرای دستورات مهاجرت ایجاد کنید:

\begin{sphinxVerbatim}[commandchars=\\\{\}]
python\PYG{+w}{ }manage.py\PYG{+w}{ }migrate
\end{sphinxVerbatim}

\item {} 
\sphinxAtStartPar
\sphinxstylestrong{ایجاد کاربر مدیر}

\sphinxAtStartPar
یک کاربر مدیر برای پنل ادمین ایجاد کنید:

\begin{sphinxVerbatim}[commandchars=\\\{\}]
python\PYG{+w}{ }manage.py\PYG{+w}{ }createsuperuser
\end{sphinxVerbatim}

\end{enumerate}


\subsection{پیکربندی}
\label{\detokenize{installation/index:id4}}
\sphinxAtStartPar
پیش از راه‌اندازی پروژه، نیاز به پیکربندی برخی تنظیمات دارید:
\begin{enumerate}
\sphinxsetlistlabels{\arabic}{enumi}{enumii}{}{.}%
\item {} 
\sphinxAtStartPar
\sphinxstylestrong{تنظیمات پایگاه داده}

\sphinxAtStartPar
در فایل settings.py، تنظیمات پایگاه داده را بر اساس محیط خود تنظیم کنید:

\begin{sphinxVerbatim}[commandchars=\\\{\}]
\PYG{n}{DATABASES} \PYG{o}{=} \PYG{p}{\PYGZob{}}
    \PYG{l+s+s1}{\PYGZsq{}}\PYG{l+s+s1}{default}\PYG{l+s+s1}{\PYGZsq{}}\PYG{p}{:} \PYG{p}{\PYGZob{}}
        \PYG{l+s+s1}{\PYGZsq{}}\PYG{l+s+s1}{ENGINE}\PYG{l+s+s1}{\PYGZsq{}}\PYG{p}{:} \PYG{l+s+s1}{\PYGZsq{}}\PYG{l+s+s1}{django.db.backends.sqlite3}\PYG{l+s+s1}{\PYGZsq{}}\PYG{p}{,}
        \PYG{l+s+s1}{\PYGZsq{}}\PYG{l+s+s1}{NAME}\PYG{l+s+s1}{\PYGZsq{}}\PYG{p}{:} \PYG{n}{BASE\PYGZus{}DIR} \PYG{o}{/} \PYG{l+s+s1}{\PYGZsq{}}\PYG{l+s+s1}{db.sqlite3}\PYG{l+s+s1}{\PYGZsq{}}\PYG{p}{,}
    \PYG{p}{\PYGZcb{}}
\PYG{p}{\PYGZcb{}}
\end{sphinxVerbatim}

\item {} 
\sphinxAtStartPar
\sphinxstylestrong{تنظیمات Redis}

\sphinxAtStartPar
برای استفاده از Celery و Django Channels، نیاز به تنظیم Redis دارید:

\begin{sphinxVerbatim}[commandchars=\\\{\}]
\PYG{n}{CELERY\PYGZus{}BROKER\PYGZus{}URL} \PYG{o}{=} \PYG{l+s+s1}{\PYGZsq{}}\PYG{l+s+s1}{redis://localhost:6379/0}\PYG{l+s+s1}{\PYGZsq{}}
\PYG{n}{CELERY\PYGZus{}RESULT\PYGZus{}BACKEND} \PYG{o}{=} \PYG{l+s+s1}{\PYGZsq{}}\PYG{l+s+s1}{redis://localhost:6379/0}\PYG{l+s+s1}{\PYGZsq{}}

\PYG{n}{CHANNEL\PYGZus{}LAYERS} \PYG{o}{=} \PYG{p}{\PYGZob{}}
    \PYG{l+s+s1}{\PYGZsq{}}\PYG{l+s+s1}{default}\PYG{l+s+s1}{\PYGZsq{}}\PYG{p}{:} \PYG{p}{\PYGZob{}}
        \PYG{l+s+s1}{\PYGZsq{}}\PYG{l+s+s1}{BACKEND}\PYG{l+s+s1}{\PYGZsq{}}\PYG{p}{:} \PYG{l+s+s1}{\PYGZsq{}}\PYG{l+s+s1}{channels\PYGZus{}redis.core.RedisChannelLayer}\PYG{l+s+s1}{\PYGZsq{}}\PYG{p}{,}
        \PYG{l+s+s1}{\PYGZsq{}}\PYG{l+s+s1}{CONFIG}\PYG{l+s+s1}{\PYGZsq{}}\PYG{p}{:} \PYG{p}{\PYGZob{}}
            \PYG{l+s+s2}{\PYGZdq{}}\PYG{l+s+s2}{hosts}\PYG{l+s+s2}{\PYGZdq{}}\PYG{p}{:} \PYG{p}{[}\PYG{p}{(}\PYG{l+s+s1}{\PYGZsq{}}\PYG{l+s+s1}{127.0.0.1}\PYG{l+s+s1}{\PYGZsq{}}\PYG{p}{,} \PYG{l+m+mi}{6379}\PYG{p}{)}\PYG{p}{]}\PYG{p}{,}
        \PYG{p}{\PYGZcb{}}\PYG{p}{,}
    \PYG{p}{\PYGZcb{}}\PYG{p}{,}
\PYG{p}{\PYGZcb{}}
\end{sphinxVerbatim}

\item {} 
\sphinxAtStartPar
\sphinxstylestrong{تنظیمات API های خارجی}

\sphinxAtStartPar
برای استفاده از API های واتساپ و پیامک، تنظیمات مربوطه را در settings.py وارد کنید:

\begin{sphinxVerbatim}[commandchars=\\\{\}]
\PYG{c+c1}{\PYGZsh{} تنظیمات واتساپ}
\PYG{n}{WHATSAPP\PYGZus{}API\PYGZus{}KEY} \PYG{o}{=} \PYG{l+s+s2}{\PYGZdq{}}\PYG{l+s+s2}{your\PYGZus{}api\PYGZus{}key}\PYG{l+s+s2}{\PYGZdq{}}
\PYG{n}{WHATSAPP\PYGZus{}API\PYGZus{}SECRET} \PYG{o}{=} \PYG{l+s+s2}{\PYGZdq{}}\PYG{l+s+s2}{your\PYGZus{}api\PYGZus{}secret}\PYG{l+s+s2}{\PYGZdq{}}

\PYG{c+c1}{\PYGZsh{} تنظیمات پیامک}
\PYG{n}{SMS\PYGZus{}API\PYGZus{}KEY} \PYG{o}{=} \PYG{l+s+s2}{\PYGZdq{}}\PYG{l+s+s2}{your\PYGZus{}sms\PYGZus{}api\PYGZus{}key}\PYG{l+s+s2}{\PYGZdq{}}
\PYG{n}{SMS\PYGZus{}API\PYGZus{}SECRET} \PYG{o}{=} \PYG{l+s+s2}{\PYGZdq{}}\PYG{l+s+s2}{your\PYGZus{}sms\PYGZus{}api\PYGZus{}secret}\PYG{l+s+s2}{\PYGZdq{}}
\end{sphinxVerbatim}

\end{enumerate}


\subsection{راه‌اندازی پروژه}
\label{\detokenize{installation/index:id5}}
\sphinxAtStartPar
پس از نصب و پیکربندی، می‌توانید پروژه را راه‌اندازی کنید:
\begin{enumerate}
\sphinxsetlistlabels{\arabic}{enumi}{enumii}{}{.}%
\item {} 
\sphinxAtStartPar
\sphinxstylestrong{اجرای سرور توسعه}

\begin{sphinxVerbatim}[commandchars=\\\{\}]
python\PYG{+w}{ }manage.py\PYG{+w}{ }runserver
\end{sphinxVerbatim}

\sphinxAtStartPar
سرور در آدرس \sphinxurl{http://localhost:8000} در دسترس خواهد بود.

\item {} 
\sphinxAtStartPar
\sphinxstylestrong{اجرای Celery}

\sphinxAtStartPar
برای اجرای وظایف غیرهمزمان، Celery را اجرا کنید:

\begin{sphinxVerbatim}[commandchars=\\\{\}]
celery\PYG{+w}{ }\PYGZhy{}A\PYG{+w}{ }config\PYG{+w}{ }worker\PYG{+w}{ }\PYGZhy{}l\PYG{+w}{ }info
\end{sphinxVerbatim}

\item {} 
\sphinxAtStartPar
\sphinxstylestrong{اجرای Celery Beat}

\sphinxAtStartPar
برای اجرای وظایف زمان‌بندی شده، Celery Beat را اجرا کنید:

\begin{sphinxVerbatim}[commandchars=\\\{\}]
celery\PYG{+w}{ }\PYGZhy{}A\PYG{+w}{ }config\PYG{+w}{ }beat\PYG{+w}{ }\PYGZhy{}l\PYG{+w}{ }info
\end{sphinxVerbatim}

\item {} 
\sphinxAtStartPar
\sphinxstylestrong{دسترسی به پنل ادمین}

\sphinxAtStartPar
با مراجعه به آدرس \sphinxurl{http://localhost:8000/ewdfsrgtyuiougyjgtrdthugghuji/} و وارد کردن نام کاربری و رمز عبور کاربر مدیر، می‌توانید به پنل ادمین دسترسی پیدا کنید.

\end{enumerate}


\subsection{عیب‌یابی متداول}
\label{\detokenize{installation/index:id6}}\begin{itemize}
\item {} 
\sphinxAtStartPar
\sphinxstylestrong{خطای اتصال به Redis}: اطمینان حاصل کنید که سرویس Redis در حال اجرا است.

\item {} 
\sphinxAtStartPar
\sphinxstylestrong{خطای اجرای مهاجرت‌ها}: اطمینان حاصل کنید که پایگاه داده به درستی پیکربندی شده است.

\item {} 
\sphinxAtStartPar
\sphinxstylestrong{خطای نصب وابستگی‌ها}: اطمینان حاصل کنید که از نسخه صحیح پایتون استفاده می‌کنید.

\end{itemize}

\sphinxstepscope


\section{معماری سیستم}
\label{\detokenize{architecture/index:architecture}}\label{\detokenize{architecture/index:id1}}\label{\detokenize{architecture/index::doc}}
\sphinxAtStartPar
در این بخش به بررسی معماری سیستم کانقالای می‌پردازیم.


\subsection{معماری کلی}
\label{\detokenize{architecture/index:id2}}
\sphinxAtStartPar
کانقالای بر اساس معماری MVC (Model\sphinxhyphen{}View\sphinxhyphen{}Controller) که در Django به صورت MTV (Model\sphinxhyphen{}Template\sphinxhyphen{}View) پیاده‌سازی شده است، طراحی شده است. این سیستم شامل چندین ماژول مستقل است که هر یک وظایف خاصی را برعهده دارند و از طریق مکانیزم‌های ارتباطی Django با یکدیگر در تعامل هستند.


\subsection{ماژول‌های اصلی}
\label{\detokenize{architecture/index:id3}}\begin{itemize}
\item {} 
\sphinxAtStartPar
\sphinxstylestrong{Core}: ماژول اصلی سیستم که مدل‌های پایه و کلاس‌های مشترک را تعریف می‌کند

\item {} 
\sphinxAtStartPar
\sphinxstylestrong{CRM}: ماژول مدیریت ارتباط با مشتری

\item {} 
\sphinxAtStartPar
\sphinxstylestrong{ERP}: ماژول مدیریت منابع سازمانی

\item {} 
\sphinxAtStartPar
\sphinxstylestrong{Sales}: ماژول مدیریت فروش

\item {} 
\sphinxAtStartPar
\sphinxstylestrong{SMS API}: ماژول ارسال و دریافت پیامک

\item {} 
\sphinxAtStartPar
\sphinxstylestrong{WhatsApp API}: ماژول ارسال و دریافت پیام در واتساپ

\item {} 
\sphinxAtStartPar
\sphinxstylestrong{Excel Manager}: ماژول مدیریت فایل‌های اکسل و گزارش‌گیری

\item {} 
\sphinxAtStartPar
\sphinxstylestrong{Chat}: ماژول سیستم گفتگوی داخلی

\end{itemize}


\subsection{فناوری‌های مورد استفاده}
\label{\detokenize{architecture/index:id4}}\begin{itemize}
\item {} 
\sphinxAtStartPar
\sphinxstylestrong{Django}: فریم‌ورک اصلی توسعه وب

\item {} 
\sphinxAtStartPar
\sphinxstylestrong{Django Channels}: برای پیاده‌سازی ارتباطات WebSocket و ASGI

\item {} 
\sphinxAtStartPar
\sphinxstylestrong{Celery}: برای مدیریت وظایف غیرهمزمان و زمان‌بندی شده

\item {} 
\sphinxAtStartPar
\sphinxstylestrong{Redis}: به عنوان broker برای Celery و سیستم کش

\item {} 
\sphinxAtStartPar
\sphinxstylestrong{SQLite}: پایگاه داده اصلی سیستم

\item {} 
\sphinxAtStartPar
\sphinxstylestrong{Django REST Framework}: برای توسعه APIs

\end{itemize}


\subsection{جریان داده}
\label{\detokenize{architecture/index:id5}}\begin{enumerate}
\sphinxsetlistlabels{\arabic}{enumi}{enumii}{}{.}%
\item {} 
\sphinxAtStartPar
درخواست‌های کاربر از طریق مرورگر یا API به سیستم ارسال می‌شود

\item {} 
\sphinxAtStartPar
URL Dispatcher درخواست را به View مناسب هدایت می‌کند

\item {} 
\sphinxAtStartPar
View با استفاده از مدل‌ها با پایگاه داده تعامل می‌کند

\item {} 
\sphinxAtStartPar
در صورت نیاز به پردازش غیرهمزمان، وظایف به Celery ارسال می‌شوند

\item {} 
\sphinxAtStartPar
در صورت نیاز به ارتباطات بلادرنگ، از Django Channels استفاده می‌شود

\item {} 
\sphinxAtStartPar
نتیجه نهایی به کاربر برگردانده می‌شود

\end{enumerate}

\sphinxstepscope


\section{ماژول‌های سیستم}
\label{\detokenize{modules/index:modules}}\label{\detokenize{modules/index:id1}}\label{\detokenize{modules/index::doc}}
\sphinxAtStartPar
در این بخش به معرفی و مستندسازی ماژول‌های مختلف سیستم کانقالای می‌پردازیم.

\sphinxstepscope


\subsection{ماژول Core}
\label{\detokenize{modules/core:core}}\label{\detokenize{modules/core:core-module}}\label{\detokenize{modules/core::doc}}
\sphinxAtStartPar
ماژول Core هسته اصلی سیستم کانقالای است که مدل‌های پایه و کلاس‌های مشترک مورد استفاده در سایر ماژول‌ها را تعریف می‌کند.


\subsubsection{مدل‌های اصلی}
\label{\detokenize{modules/core:id1}}

\paragraph{مدل Customer}
\label{\detokenize{modules/core:customer}}
\sphinxAtStartPar
مدل Customer برای نگهداری اطلاعات مشتریان طراحی شده است.

\sphinxAtStartPar
\sphinxstylestrong{فیلدهای اصلی}:
\begin{itemize}
\item {} 
\sphinxAtStartPar
\sphinxstylestrong{phone\_number}: شماره تلفن مشتری (در حال منسوخ شدن)

\item {} 
\sphinxAtStartPar
\sphinxstylestrong{sales\_agent}: ارتباط چند به چند با کارشناسان فروش

\item {} 
\sphinxAtStartPar
\sphinxstylestrong{name}: نام مشتری

\item {} 
\sphinxAtStartPar
\sphinxstylestrong{city}: شهر مشتری

\item {} 
\sphinxAtStartPar
\sphinxstylestrong{whatsapp\_preference}: ترجیح استفاده از واتساپ

\item {} 
\sphinxAtStartPar
\sphinxstylestrong{field\_of\_activity}: حیطه فعالیت مشتری

\item {} 
\sphinxAtStartPar
\sphinxstylestrong{purchase\_preferences}: ترجیحات پرداخت

\item {} 
\sphinxAtStartPar
\sphinxstylestrong{product\_preferences}: ترجیحات کالایی

\item {} 
\sphinxAtStartPar
\sphinxstylestrong{contact\_start\_time}: زمان شروع تماس

\item {} 
\sphinxAtStartPar
\sphinxstylestrong{contact\_end\_time}: زمان پایان تماس

\item {} 
\sphinxAtStartPar
\sphinxstylestrong{social\_media}: اطلاعات شبکه‌های اجتماعی

\item {} 
\sphinxAtStartPar
\sphinxstylestrong{added\_description}: توضیحات تکمیلی

\item {} 
\sphinxAtStartPar
\sphinxstylestrong{sms\_preference}: ترجیح دریافت پیامک

\item {} 
\sphinxAtStartPar
\sphinxstylestrong{day\_of\_week}: روز هفته برای ارسال پیام

\item {} 
\sphinxAtStartPar
\sphinxstylestrong{quota}: سهمیه

\item {} 
\sphinxAtStartPar
\sphinxstylestrong{created\_at}: تاریخ ثبت

\item {} 
\sphinxAtStartPar
\sphinxstylestrong{updated\_at}: تاریخ بروزرسانی

\item {} 
\sphinxAtStartPar
\sphinxstylestrong{is\_bazaargah\_enjoyer}: استفاده از بازارگاه

\end{itemize}

\sphinxAtStartPar
\sphinxstylestrong{متدهای اصلی}:
\begin{itemize}
\item {} 
\sphinxAtStartPar
\sphinxstylestrong{get\_customer\_by\_phone}: یافتن مشتری بر اساس شماره تلفن

\item {} 
\sphinxAtStartPar
\sphinxstylestrong{get\_primary\_phone}: دریافت شماره تلفن اصلی

\item {} 
\sphinxAtStartPar
\sphinxstylestrong{get\_all\_phone\_numbers}: دریافت تمام شماره‌های تلفن

\item {} 
\sphinxAtStartPar
\sphinxstylestrong{get\_tasks\_by\_time}: دریافت وظایف مرتبط با مشتری بر اساس زمان

\item {} 
\sphinxAtStartPar
\sphinxstylestrong{get\_calls\_by\_time}: دریافت تماس‌های مرتبط با مشتری بر اساس زمان

\item {} 
\sphinxAtStartPar
\sphinxstylestrong{get\_purchases\_by\_time}: دریافت خریدهای مشتری بر اساس زمان

\item {} 
\sphinxAtStartPar
\sphinxstylestrong{get\_field\_of\_activity\_display}: نمایش حیطه فعالیت مشتری

\item {} 
\sphinxAtStartPar
\sphinxstylestrong{get\_field\_of\_activity\_list}: دریافت لیست حیطه‌های فعالیت

\item {} 
\sphinxAtStartPar
\sphinxstylestrong{set\_field\_of\_activity\_list}: تنظیم لیست حیطه‌های فعالیت

\end{itemize}


\paragraph{مدل SalesAgent}
\label{\detokenize{modules/core:salesagent}}
\sphinxAtStartPar
مدل SalesAgent برای نگهداری اطلاعات کارشناسان فروش طراحی شده است.

\sphinxAtStartPar
\sphinxstylestrong{فیلدهای اصلی}:
\begin{itemize}
\item {} 
\sphinxAtStartPar
\sphinxstylestrong{user}: ارتباط یک به یک با مدل User

\item {} 
\sphinxAtStartPar
\sphinxstylestrong{name}: نام کارشناس فروش

\item {} 
\sphinxAtStartPar
\sphinxstylestrong{email}: ایمیل کارشناس فروش

\item {} 
\sphinxAtStartPar
\sphinxstylestrong{phone}: شماره تلفن کارشناس فروش

\item {} 
\sphinxAtStartPar
\sphinxstylestrong{created\_at}: تاریخ ثبت

\item {} 
\sphinxAtStartPar
\sphinxstylestrong{updated\_at}: تاریخ بروزرسانی

\end{itemize}

\sphinxAtStartPar
\sphinxstylestrong{متدهای اصلی}:
\begin{itemize}
\item {} 
\sphinxAtStartPar
\sphinxstylestrong{get\_customers}: دریافت مشتریان مرتبط با کارشناس فروش

\item {} 
\sphinxAtStartPar
\sphinxstylestrong{get\_tasks}: دریافت وظایف مرتبط با کارشناس فروش

\item {} 
\sphinxAtStartPar
\sphinxstylestrong{get\_calls}: دریافت تماس‌های مرتبط با کارشناس فروش

\item {} 
\sphinxAtStartPar
\sphinxstylestrong{get\_purchases}: دریافت خریدهای مرتبط با کارشناس فروش

\end{itemize}


\paragraph{مدل CargoProduct}
\label{\detokenize{modules/core:cargoproduct}}
\sphinxAtStartPar
مدل CargoProduct برای نگهداری اطلاعات محصولات طراحی شده است.

\sphinxAtStartPar
\sphinxstylestrong{فیلدهای اصلی}:
\begin{itemize}
\item {} 
\sphinxAtStartPar
\sphinxstylestrong{product\_name}: نام محصول

\item {} 
\sphinxAtStartPar
\sphinxstylestrong{product\_type}: نوع محصول

\item {} 
\sphinxAtStartPar
\sphinxstylestrong{brand}: برند محصول

\item {} 
\sphinxAtStartPar
\sphinxstylestrong{model}: مدل محصول

\item {} 
\sphinxAtStartPar
\sphinxstylestrong{origin}: مبدا محصول

\item {} 
\sphinxAtStartPar
\sphinxstylestrong{volume}: حجم محصول

\item {} 
\sphinxAtStartPar
\sphinxstylestrong{weight}: وزن محصول

\item {} 
\sphinxAtStartPar
\sphinxstylestrong{description}: توضیحات محصول

\item {} 
\sphinxAtStartPar
\sphinxstylestrong{barcode}: بارکد محصول

\item {} 
\sphinxAtStartPar
\sphinxstylestrong{created\_at}: تاریخ ثبت

\item {} 
\sphinxAtStartPar
\sphinxstylestrong{updated\_at}: تاریخ بروزرسانی

\end{itemize}


\paragraph{مدل PurchaseModel}
\label{\detokenize{modules/core:purchasemodel}}
\sphinxAtStartPar
مدل PurchaseModel برای نگهداری اطلاعات خریدها طراحی شده است.

\sphinxAtStartPar
\sphinxstylestrong{فیلدهای اصلی}:
\begin{itemize}
\item {} 
\sphinxAtStartPar
\sphinxstylestrong{customer}: ارتباط با مشتری

\item {} 
\sphinxAtStartPar
\sphinxstylestrong{sales\_agent}: ارتباط با کارشناس فروش

\item {} 
\sphinxAtStartPar
\sphinxstylestrong{product}: ارتباط با محصول

\item {} 
\sphinxAtStartPar
\sphinxstylestrong{quantity}: تعداد

\item {} 
\sphinxAtStartPar
\sphinxstylestrong{price}: قیمت

\item {} 
\sphinxAtStartPar
\sphinxstylestrong{total\_price}: قیمت کل

\item {} 
\sphinxAtStartPar
\sphinxstylestrong{purchase\_date}: تاریخ خرید

\item {} 
\sphinxAtStartPar
\sphinxstylestrong{status}: وضعیت خرید

\item {} 
\sphinxAtStartPar
\sphinxstylestrong{payment\_method}: روش پرداخت

\item {} 
\sphinxAtStartPar
\sphinxstylestrong{created\_at}: تاریخ ثبت

\item {} 
\sphinxAtStartPar
\sphinxstylestrong{updated\_at}: تاریخ بروزرسانی

\end{itemize}


\paragraph{مدل TaskModel}
\label{\detokenize{modules/core:taskmodel}}
\sphinxAtStartPar
مدل TaskModel برای نگهداری اطلاعات وظایف طراحی شده است.

\sphinxAtStartPar
\sphinxstylestrong{فیلدهای اصلی}:
\begin{itemize}
\item {} 
\sphinxAtStartPar
\sphinxstylestrong{title}: عنوان وظیفه

\item {} 
\sphinxAtStartPar
\sphinxstylestrong{description}: توضیحات وظیفه

\item {} 
\sphinxAtStartPar
\sphinxstylestrong{assigned\_to}: ارتباط با کاربر مسئول

\item {} 
\sphinxAtStartPar
\sphinxstylestrong{customer}: ارتباط با مشتری

\item {} 
\sphinxAtStartPar
\sphinxstylestrong{due\_date}: تاریخ سررسید

\item {} 
\sphinxAtStartPar
\sphinxstylestrong{status}: وضعیت وظیفه

\item {} 
\sphinxAtStartPar
\sphinxstylestrong{priority}: اولویت وظیفه

\item {} 
\sphinxAtStartPar
\sphinxstylestrong{created\_at}: تاریخ ثبت

\item {} 
\sphinxAtStartPar
\sphinxstylestrong{updated\_at}: تاریخ بروزرسانی

\end{itemize}


\paragraph{مدل NotificationModel}
\label{\detokenize{modules/core:notificationmodel}}
\sphinxAtStartPar
مدل NotificationModel برای نگهداری اطلاعات اعلان‌ها طراحی شده است.

\sphinxAtStartPar
\sphinxstylestrong{فیلدهای اصلی}:
\begin{itemize}
\item {} 
\sphinxAtStartPar
\sphinxstylestrong{user}: ارتباط با کاربر

\item {} 
\sphinxAtStartPar
\sphinxstylestrong{title}: عنوان اعلان

\item {} 
\sphinxAtStartPar
\sphinxstylestrong{message}: متن اعلان

\item {} 
\sphinxAtStartPar
\sphinxstylestrong{type}: نوع اعلان

\item {} 
\sphinxAtStartPar
\sphinxstylestrong{read}: خوانده شده/نشده

\item {} 
\sphinxAtStartPar
\sphinxstylestrong{created\_at}: تاریخ ثبت

\end{itemize}


\subsubsection{فرم‌ها}
\label{\detokenize{modules/core:id2}}
\sphinxAtStartPar
در ماژول Core، فرم‌های مختلفی برای کار با مدل‌های فوق تعریف شده‌اند، از جمله:
\begin{itemize}
\item {} 
\sphinxAtStartPar
\sphinxstylestrong{CustomerForm}: فرم ایجاد و ویرایش مشتری

\item {} 
\sphinxAtStartPar
\sphinxstylestrong{SalesAgentForm}: فرم ایجاد و ویرایش کارشناس فروش

\item {} 
\sphinxAtStartPar
\sphinxstylestrong{ProductForm}: فرم ایجاد و ویرایش محصول

\item {} 
\sphinxAtStartPar
\sphinxstylestrong{PurchaseForm}: فرم ایجاد و ویرایش خرید

\item {} 
\sphinxAtStartPar
\sphinxstylestrong{TaskForm}: فرم ایجاد و ویرایش وظیفه

\end{itemize}


\subsubsection{میان‌افزارها}
\label{\detokenize{modules/core:id3}}
\sphinxAtStartPar
ماژول Core شامل چندین میان‌افزار (Middleware) است که وظایف مختلفی را انجام می‌دهند:
\begin{itemize}
\item {} 
\sphinxAtStartPar
\sphinxstylestrong{SessionTimeoutMiddleware}: میان‌افزار مدیریت زمان‌سنجی نشست کاربر

\item {} 
\sphinxAtStartPar
\sphinxstylestrong{AuthenticationMiddleware}: میان‌افزار احراز هویت کاربر

\end{itemize}


\subsubsection{دکوراتورها}
\label{\detokenize{modules/core:id4}}
\sphinxAtStartPar
ماژول Core شامل چندین دکوراتور است که برای مدیریت دسترسی و کنترل رفتار توابع view طراحی شده‌اند:
\begin{itemize}
\item {} 
\sphinxAtStartPar
\sphinxstylestrong{login\_required}: بررسی لاگین بودن کاربر

\item {} 
\sphinxAtStartPar
\sphinxstylestrong{permission\_required}: بررسی دسترسی‌های کاربر

\item {} 
\sphinxAtStartPar
\sphinxstylestrong{sales\_agent\_required}: بررسی کارشناس فروش بودن کاربر

\end{itemize}

\sphinxstepscope


\subsection{ماژول WhatsApp API}
\label{\detokenize{modules/whatsapp_api:whatsapp-api}}\label{\detokenize{modules/whatsapp_api:whatsapp-api-module}}\label{\detokenize{modules/whatsapp_api::doc}}
\sphinxAtStartPar
ماژول WhatsApp API برای ارسال و دریافت پیام‌های واتساپ طراحی شده است. این ماژول با API واتساپ یکپارچه‌سازی شده و امکان ارسال پیام به مشتریان و دریافت پاسخ آن‌ها را فراهم می‌کند.


\subsubsection{مدل‌های اصلی}
\label{\detokenize{modules/whatsapp_api:id1}}

\paragraph{مدل WhatsAppMessage}
\label{\detokenize{modules/whatsapp_api:whatsappmessage}}
\sphinxAtStartPar
مدل WhatsAppMessage برای نگهداری اطلاعات پیام‌های واتساپ طراحی شده است.

\sphinxAtStartPar
\sphinxstylestrong{فیلدهای اصلی}:
\begin{itemize}
\item {} 
\sphinxAtStartPar
\sphinxstylestrong{sender}: فرستنده پیام

\item {} 
\sphinxAtStartPar
\sphinxstylestrong{recipient}: گیرنده پیام

\item {} 
\sphinxAtStartPar
\sphinxstylestrong{message}: متن پیام

\item {} 
\sphinxAtStartPar
\sphinxstylestrong{media\_url}: آدرس محتوای چندرسانه‌ای

\item {} 
\sphinxAtStartPar
\sphinxstylestrong{media\_type}: نوع محتوای چندرسانه‌ای

\item {} 
\sphinxAtStartPar
\sphinxstylestrong{status}: وضعیت پیام

\item {} 
\sphinxAtStartPar
\sphinxstylestrong{message\_id}: شناسه پیام در واتساپ

\item {} 
\sphinxAtStartPar
\sphinxstylestrong{customer}: ارتباط با مشتری

\item {} 
\sphinxAtStartPar
\sphinxstylestrong{sales\_agent}: ارتباط با کارشناس فروش

\item {} 
\sphinxAtStartPar
\sphinxstylestrong{sent\_at}: تاریخ ارسال

\item {} 
\sphinxAtStartPar
\sphinxstylestrong{received\_at}: تاریخ دریافت

\item {} 
\sphinxAtStartPar
\sphinxstylestrong{created\_at}: تاریخ ثبت

\item {} 
\sphinxAtStartPar
\sphinxstylestrong{updated\_at}: تاریخ بروزرسانی

\end{itemize}


\paragraph{مدل WhatsAppTemplate}
\label{\detokenize{modules/whatsapp_api:whatsapptemplate}}
\sphinxAtStartPar
مدل WhatsAppTemplate برای نگهداری اطلاعات قالب‌های پیام واتساپ طراحی شده است.

\sphinxAtStartPar
\sphinxstylestrong{فیلدهای اصلی}:
\begin{itemize}
\item {} 
\sphinxAtStartPar
\sphinxstylestrong{name}: نام قالب

\item {} 
\sphinxAtStartPar
\sphinxstylestrong{type}: نوع قالب

\item {} 
\sphinxAtStartPar
\sphinxstylestrong{category}: دسته‌بندی قالب

\item {} 
\sphinxAtStartPar
\sphinxstylestrong{components}: اجزای قالب

\item {} 
\sphinxAtStartPar
\sphinxstylestrong{language}: زبان قالب

\item {} 
\sphinxAtStartPar
\sphinxstylestrong{status}: وضعیت قالب

\item {} 
\sphinxAtStartPar
\sphinxstylestrong{created\_at}: تاریخ ثبت

\item {} 
\sphinxAtStartPar
\sphinxstylestrong{updated\_at}: تاریخ بروزرسانی

\end{itemize}


\paragraph{مدل WhatsAppAccount}
\label{\detokenize{modules/whatsapp_api:whatsappaccount}}
\sphinxAtStartPar
مدل WhatsAppAccount برای نگهداری اطلاعات حساب‌های واتساپ طراحی شده است.

\sphinxAtStartPar
\sphinxstylestrong{فیلدهای اصلی}:
\begin{itemize}
\item {} 
\sphinxAtStartPar
\sphinxstylestrong{phone\_number}: شماره تلفن

\item {} 
\sphinxAtStartPar
\sphinxstylestrong{business\_name}: نام کسب و کار

\item {} 
\sphinxAtStartPar
\sphinxstylestrong{api\_key}: کلید API

\item {} 
\sphinxAtStartPar
\sphinxstylestrong{api\_secret}: رمز API

\item {} 
\sphinxAtStartPar
\sphinxstylestrong{webhook\_url}: آدرس وب‌هوک

\item {} 
\sphinxAtStartPar
\sphinxstylestrong{status}: وضعیت حساب

\item {} 
\sphinxAtStartPar
\sphinxstylestrong{created\_at}: تاریخ ثبت

\item {} 
\sphinxAtStartPar
\sphinxstylestrong{updated\_at}: تاریخ بروزرسانی

\end{itemize}


\subsubsection{توابع اصلی}
\label{\detokenize{modules/whatsapp_api:id2}}

\paragraph{تابع send\_whatsapp\_message}
\label{\detokenize{modules/whatsapp_api:send-whatsapp-message}}
\sphinxAtStartPar
این تابع برای ارسال پیام واتساپ استفاده می‌شود.

\sphinxAtStartPar
\sphinxstylestrong{پارامترها}:
\begin{itemize}
\item {} 
\sphinxAtStartPar
\sphinxstylestrong{to}: شماره تلفن گیرنده

\item {} 
\sphinxAtStartPar
\sphinxstylestrong{message}: متن پیام

\item {} 
\sphinxAtStartPar
\sphinxstylestrong{media\_url}: آدرس محتوای چندرسانه‌ای (اختیاری)

\item {} 
\sphinxAtStartPar
\sphinxstylestrong{media\_type}: نوع محتوای چندرسانه‌ای (اختیاری)

\item {} 
\sphinxAtStartPar
\sphinxstylestrong{customer\_id}: شناسه مشتری (اختیاری)

\item {} 
\sphinxAtStartPar
\sphinxstylestrong{sales\_agent\_id}: شناسه کارشناس فروش (اختیاری)

\end{itemize}

\sphinxAtStartPar
\sphinxstylestrong{مثال استفاده}:

\begin{sphinxVerbatim}[commandchars=\\\{\}]
\PYG{k+kn}{from}\PYG{+w}{ }\PYG{n+nn}{whatsapp\PYGZus{}api}\PYG{n+nn}{.}\PYG{n+nn}{utils}\PYG{+w}{ }\PYG{k+kn}{import} \PYG{n}{send\PYGZus{}whatsapp\PYGZus{}message}

\PYG{n}{send\PYGZus{}whatsapp\PYGZus{}message}\PYG{p}{(}
    \PYG{n}{to}\PYG{o}{=}\PYG{l+s+s2}{\PYGZdq{}}\PYG{l+s+s2}{989123456789}\PYG{l+s+s2}{\PYGZdq{}}\PYG{p}{,}
    \PYG{n}{message}\PYG{o}{=}\PYG{l+s+s2}{\PYGZdq{}}\PYG{l+s+s2}{سلام، خوش آمدید!}\PYG{l+s+s2}{\PYGZdq{}}\PYG{p}{,}
    \PYG{n}{customer\PYGZus{}id}\PYG{o}{=}\PYG{l+m+mi}{1}
\PYG{p}{)}
\end{sphinxVerbatim}


\paragraph{تابع send\_whatsapp\_template}
\label{\detokenize{modules/whatsapp_api:send-whatsapp-template}}
\sphinxAtStartPar
این تابع برای ارسال پیام واتساپ با استفاده از قالب استفاده می‌شود.

\sphinxAtStartPar
\sphinxstylestrong{پارامترها}:
\begin{itemize}
\item {} 
\sphinxAtStartPar
\sphinxstylestrong{to}: شماره تلفن گیرنده

\item {} 
\sphinxAtStartPar
\sphinxstylestrong{template\_name}: نام قالب

\item {} 
\sphinxAtStartPar
\sphinxstylestrong{template\_params}: پارامترهای قالب

\item {} 
\sphinxAtStartPar
\sphinxstylestrong{customer\_id}: شناسه مشتری (اختیاری)

\item {} 
\sphinxAtStartPar
\sphinxstylestrong{sales\_agent\_id}: شناسه کارشناس فروش (اختیاری)

\end{itemize}

\sphinxAtStartPar
\sphinxstylestrong{مثال استفاده}:

\begin{sphinxVerbatim}[commandchars=\\\{\}]
\PYG{k+kn}{from}\PYG{+w}{ }\PYG{n+nn}{whatsapp\PYGZus{}api}\PYG{n+nn}{.}\PYG{n+nn}{utils}\PYG{+w}{ }\PYG{k+kn}{import} \PYG{n}{send\PYGZus{}whatsapp\PYGZus{}template}

\PYG{n}{send\PYGZus{}whatsapp\PYGZus{}template}\PYG{p}{(}
    \PYG{n}{to}\PYG{o}{=}\PYG{l+s+s2}{\PYGZdq{}}\PYG{l+s+s2}{989123456789}\PYG{l+s+s2}{\PYGZdq{}}\PYG{p}{,}
    \PYG{n}{template\PYGZus{}name}\PYG{o}{=}\PYG{l+s+s2}{\PYGZdq{}}\PYG{l+s+s2}{welcome\PYGZus{}message}\PYG{l+s+s2}{\PYGZdq{}}\PYG{p}{,}
    \PYG{n}{template\PYGZus{}params}\PYG{o}{=}\PYG{p}{\PYGZob{}}\PYG{l+s+s2}{\PYGZdq{}}\PYG{l+s+s2}{name}\PYG{l+s+s2}{\PYGZdq{}}\PYG{p}{:} \PYG{l+s+s2}{\PYGZdq{}}\PYG{l+s+s2}{علی}\PYG{l+s+s2}{\PYGZdq{}}\PYG{p}{,} \PYG{l+s+s2}{\PYGZdq{}}\PYG{l+s+s2}{company}\PYG{l+s+s2}{\PYGZdq{}}\PYG{p}{:} \PYG{l+s+s2}{\PYGZdq{}}\PYG{l+s+s2}{کانقالای}\PYG{l+s+s2}{\PYGZdq{}}\PYG{p}{\PYGZcb{}}\PYG{p}{,}
    \PYG{n}{customer\PYGZus{}id}\PYG{o}{=}\PYG{l+m+mi}{1}
\PYG{p}{)}
\end{sphinxVerbatim}


\subsubsection{وظایف Celery}
\label{\detokenize{modules/whatsapp_api:celery}}

\paragraph{وظیفه send\_bulk\_whatsapp\_messages}
\label{\detokenize{modules/whatsapp_api:send-bulk-whatsapp-messages}}
\sphinxAtStartPar
این وظیفه برای ارسال پیام واتساپ به صورت انبوه استفاده می‌شود.

\sphinxAtStartPar
\sphinxstylestrong{پارامترها}:
\begin{itemize}
\item {} 
\sphinxAtStartPar
\sphinxstylestrong{numbers}: لیست شماره تلفن‌ها

\item {} 
\sphinxAtStartPar
\sphinxstylestrong{message}: متن پیام

\item {} 
\sphinxAtStartPar
\sphinxstylestrong{media\_url}: آدرس محتوای چندرسانه‌ای (اختیاری)

\item {} 
\sphinxAtStartPar
\sphinxstylestrong{media\_type}: نوع محتوای چندرسانه‌ای (اختیاری)

\end{itemize}

\sphinxAtStartPar
\sphinxstylestrong{مثال استفاده}:

\begin{sphinxVerbatim}[commandchars=\\\{\}]
\PYG{k+kn}{from}\PYG{+w}{ }\PYG{n+nn}{whatsapp\PYGZus{}api}\PYG{n+nn}{.}\PYG{n+nn}{tasks}\PYG{+w}{ }\PYG{k+kn}{import} \PYG{n}{send\PYGZus{}bulk\PYGZus{}whatsapp\PYGZus{}messages}

\PYG{n}{send\PYGZus{}bulk\PYGZus{}whatsapp\PYGZus{}messages}\PYG{o}{.}\PYG{n}{delay}\PYG{p}{(}
    \PYG{n}{numbers}\PYG{o}{=}\PYG{p}{[}\PYG{l+s+s2}{\PYGZdq{}}\PYG{l+s+s2}{989123456789}\PYG{l+s+s2}{\PYGZdq{}}\PYG{p}{,} \PYG{l+s+s2}{\PYGZdq{}}\PYG{l+s+s2}{989123456788}\PYG{l+s+s2}{\PYGZdq{}}\PYG{p}{]}\PYG{p}{,}
    \PYG{n}{message}\PYG{o}{=}\PYG{l+s+s2}{\PYGZdq{}}\PYG{l+s+s2}{سلام، خوش آمدید!}\PYG{l+s+s2}{\PYGZdq{}}
\PYG{p}{)}
\end{sphinxVerbatim}


\paragraph{وظیفه process\_whatsapp\_webhook}
\label{\detokenize{modules/whatsapp_api:process-whatsapp-webhook}}
\sphinxAtStartPar
این وظیفه برای پردازش وب‌هوک‌های واتساپ استفاده می‌شود.

\sphinxAtStartPar
\sphinxstylestrong{پارامترها}:
\begin{itemize}
\item {} 
\sphinxAtStartPar
\sphinxstylestrong{data}: داده‌های وب‌هوک

\end{itemize}

\sphinxAtStartPar
\sphinxstylestrong{مثال استفاده}:

\begin{sphinxVerbatim}[commandchars=\\\{\}]
\PYG{k+kn}{from}\PYG{+w}{ }\PYG{n+nn}{whatsapp\PYGZus{}api}\PYG{n+nn}{.}\PYG{n+nn}{tasks}\PYG{+w}{ }\PYG{k+kn}{import} \PYG{n}{process\PYGZus{}whatsapp\PYGZus{}webhook}

\PYG{n}{process\PYGZus{}whatsapp\PYGZus{}webhook}\PYG{o}{.}\PYG{n}{delay}\PYG{p}{(}\PYG{n}{data}\PYG{o}{=}\PYG{n}{webhook\PYGZus{}data}\PYG{p}{)}
\end{sphinxVerbatim}


\subsubsection{تنظیم و پیکربندی}
\label{\detokenize{modules/whatsapp_api:id3}}
\sphinxAtStartPar
برای استفاده از ماژول WhatsApp API، نیاز به تنظیم موارد زیر است:
\begin{enumerate}
\sphinxsetlistlabels{\arabic}{enumi}{enumii}{}{.}%
\item {} 
\sphinxAtStartPar
ثبت‌نام در سرویس واتساپ بیزینس API

\item {} 
\sphinxAtStartPar
دریافت کلید API

\item {} 
\sphinxAtStartPar
تنظیم وب‌هوک

\item {} 
\sphinxAtStartPar
تنظیم قالب‌های پیام

\end{enumerate}

\sphinxAtStartPar
\sphinxstylestrong{مثال تنظیمات در فایل settings.py}:

\begin{sphinxVerbatim}[commandchars=\\\{\}]
\PYG{c+c1}{\PYGZsh{} تنظیمات واتساپ}
\PYG{n}{WHATSAPP\PYGZus{}API\PYGZus{}KEY} \PYG{o}{=} \PYG{l+s+s2}{\PYGZdq{}}\PYG{l+s+s2}{your\PYGZus{}api\PYGZus{}key}\PYG{l+s+s2}{\PYGZdq{}}
\PYG{n}{WHATSAPP\PYGZus{}API\PYGZus{}SECRET} \PYG{o}{=} \PYG{l+s+s2}{\PYGZdq{}}\PYG{l+s+s2}{your\PYGZus{}api\PYGZus{}secret}\PYG{l+s+s2}{\PYGZdq{}}
\PYG{n}{WHATSAPP\PYGZus{}PHONE\PYGZus{}NUMBER} \PYG{o}{=} \PYG{l+s+s2}{\PYGZdq{}}\PYG{l+s+s2}{your\PYGZus{}phone\PYGZus{}number}\PYG{l+s+s2}{\PYGZdq{}}
\PYG{n}{WHATSAPP\PYGZus{}WEBHOOK\PYGZus{}URL} \PYG{o}{=} \PYG{l+s+s2}{\PYGZdq{}}\PYG{l+s+s2}{your\PYGZus{}webhook\PYGZus{}url}\PYG{l+s+s2}{\PYGZdq{}}
\end{sphinxVerbatim}


\subsubsection{API Endpoints}
\label{\detokenize{modules/whatsapp_api:api-endpoints}}
\sphinxAtStartPar
ماژول WhatsApp API چندین API endpoint برای ارسال و دریافت پیام‌های واتساپ فراهم می‌کند:
\begin{itemize}
\item {} 
\sphinxAtStartPar
\sphinxstylestrong{/whatsapp/send/}: ارسال پیام واتساپ

\item {} 
\sphinxAtStartPar
\sphinxstylestrong{/whatsapp/send\sphinxhyphen{}template/}: ارسال پیام واتساپ با استفاده از قالب

\item {} 
\sphinxAtStartPar
\sphinxstylestrong{/whatsapp/webhook/}: دریافت وب‌هوک‌های واتساپ

\item {} 
\sphinxAtStartPar
\sphinxstylestrong{/whatsapp/messages/}: دریافت لیست پیام‌های واتساپ

\item {} 
\sphinxAtStartPar
\sphinxstylestrong{/whatsapp/templates/}: دریافت لیست قالب‌های واتساپ

\end{itemize}


\subsubsection{داشبورد}
\label{\detokenize{modules/whatsapp_api:id4}}
\sphinxAtStartPar
ماژول WhatsApp API دارای یک داشبورد برای مدیریت پیام‌های واتساپ است. این داشبورد در آدرس زیر قابل دسترسی است:

\begin{sphinxVerbatim}[commandchars=\\\{\}]
/whatsapp/dashboard/
\end{sphinxVerbatim}

\sphinxAtStartPar
در این داشبورد، امکانات زیر فراهم شده است:
\begin{itemize}
\item {} 
\sphinxAtStartPar
مشاهده لیست پیام‌های ارسالی و دریافتی

\item {} 
\sphinxAtStartPar
ارسال پیام جدید

\item {} 
\sphinxAtStartPar
ارسال پیام انبوه

\item {} 
\sphinxAtStartPar
مدیریت قالب‌های پیام

\item {} 
\sphinxAtStartPar
مشاهده گزارش‌های ارسال پیام

\end{itemize}


\subsection{ماژول Core}
\label{\detokenize{modules/index:core}}
\sphinxAtStartPar
ماژول Core هسته اصلی سیستم است که مدل‌های پایه و کلاس‌های مشترک را تعریف می‌کند. این ماژول شامل مدل‌هایی مانند مشتری، کارشناس فروش، محصول و غیره است.

\sphinxAtStartPar
مدل‌های اصلی:
\begin{itemize}
\item {} 
\sphinxAtStartPar
\sphinxstylestrong{Customer}: مدل مشتری

\item {} 
\sphinxAtStartPar
\sphinxstylestrong{SalesAgent}: مدل کارشناس فروش

\item {} 
\sphinxAtStartPar
\sphinxstylestrong{CargoProduct}: مدل محصول

\item {} 
\sphinxAtStartPar
\sphinxstylestrong{PurchaseModel}: مدل خرید

\item {} 
\sphinxAtStartPar
\sphinxstylestrong{TaskModel}: مدل وظیفه

\item {} 
\sphinxAtStartPar
\sphinxstylestrong{NotificationModel}: مدل اعلان

\end{itemize}


\subsection{ماژول CRM}
\label{\detokenize{modules/index:crm}}
\sphinxAtStartPar
ماژول CRM (Customer Relationship Management) برای مدیریت ارتباط با مشتری طراحی شده است. این ماژول امکاناتی مانند ثبت و پیگیری تماس‌ها، مدیریت مشتریان، و ثبت وظایف مرتبط با مشتریان را فراهم می‌کند.


\subsection{ماژول ERP}
\label{\detokenize{modules/index:erp}}
\sphinxAtStartPar
ماژول ERP (Enterprise Resource Planning) برای مدیریت منابع سازمانی طراحی شده است. این ماژول امکاناتی مانند مدیریت محصولات، مدیریت موجودی، و مدیریت فرآیندهای داخلی سازمان را فراهم می‌کند.


\subsection{ماژول Sales}
\label{\detokenize{modules/index:sales}}
\sphinxAtStartPar
ماژول Sales برای مدیریت فروش طراحی شده است. این ماژول امکاناتی مانند ثبت سفارش، پیگیری وضعیت سفارش، و گزارش‌گیری از فروش را فراهم می‌کند.


\subsection{ماژول SMS API}
\label{\detokenize{modules/index:sms-api}}
\sphinxAtStartPar
ماژول SMS API برای ارسال و دریافت پیامک طراحی شده است. این ماژول با سرویس‌های پیامکی مختلف یکپارچه‌سازی شده و امکان ارسال پیامک به مشتریان و دریافت پاسخ آن‌ها را فراهم می‌کند.


\subsection{ماژول WhatsApp API}
\label{\detokenize{modules/index:whatsapp-api}}
\sphinxAtStartPar
ماژول WhatsApp API برای ارسال و دریافت پیام در واتساپ طراحی شده است. این ماژول با API واتساپ یکپارچه‌سازی شده و امکان ارسال پیام به مشتریان و دریافت پاسخ آن‌ها را فراهم می‌کند.


\subsection{ماژول Excel Manager}
\label{\detokenize{modules/index:excel-manager}}
\sphinxAtStartPar
ماژول Excel Manager برای مدیریت فایل‌های اکسل و گزارش‌گیری طراحی شده است. این ماژول امکاناتی مانند وارد کردن داده از فایل اکسل، تولید گزارش‌های اکسل، و تجزیه و تحلیل داده‌ها را فراهم می‌کند.


\subsection{ماژول Chat}
\label{\detokenize{modules/index:chat}}
\sphinxAtStartPar
ماژول Chat برای ارتباطات داخلی سیستم طراحی شده است. این ماژول با استفاده از Django Channels و WebSockets پیاده‌سازی شده و امکان گفتگوی بلادرنگ بین کاربران سیستم را فراهم می‌کند.

\sphinxstepscope


\section{راهنمای کاربری}
\label{\detokenize{user_guide/index:user-guide}}\label{\detokenize{user_guide/index:id1}}\label{\detokenize{user_guide/index::doc}}
\sphinxAtStartPar
در این بخش، راهنمای کاربری سیستم کانقالای ارائه شده است. این راهنما به کاربران کمک می‌کند تا با قابلیت‌های مختلف سیستم آشنا شوند و بتوانند از آن به بهترین شکل استفاده کنند.


\subsection{مقدمه}
\label{\detokenize{user_guide/index:id2}}
\sphinxAtStartPar
کانقالای یک سیستم مدیریت منابع سازمانی (ERP) و مدیریت ارتباط با مشتری (CRM) است که برای مدیریت عملیات کسب و کار طراحی شده است. این سیستم امکان مدیریت مشتریان، فروش، محصولات، و ارتباطات را فراهم می‌کند.


\subsection{ورود به سیستم}
\label{\detokenize{user_guide/index:id3}}
\sphinxAtStartPar
برای ورود به سیستم، به آدرس وب سایت کانقالای مراجعه کنید و از صفحه ورود، اطلاعات حساب کاربری خود را وارد کنید:
\begin{enumerate}
\sphinxsetlistlabels{\arabic}{enumi}{enumii}{}{.}%
\item {} 
\sphinxAtStartPar
نام کاربری و رمز عبور خود را وارد کنید

\item {} 
\sphinxAtStartPar
دکمه "ورود" را کلیک کنید

\item {} 
\sphinxAtStartPar
پس از احراز هویت، به داشبورد اصلی منتقل خواهید شد

\end{enumerate}

\noindent{\hspace*{\fill}\sphinxincludegraphics{{user_guide\_static\images\login}.png}\hspace*{\fill}}


\subsection{داشبورد}
\label{\detokenize{user_guide/index:id4}}
\sphinxAtStartPar
داشبورد اصلی، صفحه اصلی سیستم است که پس از ورود به سیستم نمایش داده می‌شود. در این صفحه، اطلاعات کلی و آمار مهم نمایش داده می‌شود:
\begin{itemize}
\item {} 
\sphinxAtStartPar
\sphinxstylestrong{خلاصه آمار}: تعداد مشتریان، فروش‌ها، وظایف و غیره

\item {} 
\sphinxAtStartPar
\sphinxstylestrong{نمودارها}: نمودارهای آماری مانند فروش ماهانه، تعداد مشتریان جدید و غیره

\item {} 
\sphinxAtStartPar
\sphinxstylestrong{وظایف اخیر}: لیست وظایف اخیر کاربر

\item {} 
\sphinxAtStartPar
\sphinxstylestrong{اعلان‌ها}: اعلان‌های مهم سیستم

\end{itemize}


\subsection{منوی اصلی}
\label{\detokenize{user_guide/index:id5}}
\sphinxAtStartPar
منوی اصلی سیستم در سمت راست صفحه قرار دارد و دسترسی به بخش‌های مختلف سیستم را فراهم می‌کند:
\begin{itemize}
\item {} 
\sphinxAtStartPar
\sphinxstylestrong{داشبورد}: صفحه اصلی سیستم

\item {} 
\sphinxAtStartPar
\sphinxstylestrong{مشتریان}: مدیریت مشتریان

\item {} 
\sphinxAtStartPar
\sphinxstylestrong{فروش}: مدیریت فروش

\item {} 
\sphinxAtStartPar
\sphinxstylestrong{محصولات}: مدیریت محصولات

\item {} 
\sphinxAtStartPar
\sphinxstylestrong{واتساپ}: مدیریت پیام‌های واتساپ

\item {} 
\sphinxAtStartPar
\sphinxstylestrong{پیامک}: مدیریت پیامک‌ها

\item {} 
\sphinxAtStartPar
\sphinxstylestrong{گزارش‌ها}: گزارش‌گیری

\item {} 
\sphinxAtStartPar
\sphinxstylestrong{تنظیمات}: تنظیمات سیستم

\item {} 
\sphinxAtStartPar
\sphinxstylestrong{مدیریت}: پنل مدیریت (فقط برای کاربران مدیر)

\end{itemize}


\subsection{مدیریت مشتریان}
\label{\detokenize{user_guide/index:id6}}
\sphinxAtStartPar
در بخش مدیریت مشتریان، می‌توانید مشتریان را مدیریت کنید:
\begin{itemize}
\item {} 
\sphinxAtStartPar
\sphinxstylestrong{افزودن مشتری جدید}: با کلیک بر روی دکمه "افزودن مشتری جدید"، می‌توانید مشتری جدید ایجاد کنید

\item {} 
\sphinxAtStartPar
\sphinxstylestrong{جستجوی مشتری}: با استفاده از فیلد جستجو، می‌توانید مشتریان را جستجو کنید

\item {} 
\sphinxAtStartPar
\sphinxstylestrong{ویرایش مشتری}: با کلیک بر روی دکمه "ویرایش" در کنار هر مشتری، می‌توانید اطلاعات مشتری را ویرایش کنید

\item {} 
\sphinxAtStartPar
\sphinxstylestrong{حذف مشتری}: با کلیک بر روی دکمه "حذف" در کنار هر مشتری، می‌توانید مشتری را حذف کنید

\item {} 
\sphinxAtStartPar
\sphinxstylestrong{مشاهده جزئیات مشتری}: با کلیک بر روی نام مشتری، می‌توانید جزئیات مشتری را مشاهده کنید

\end{itemize}


\subsection{مدیریت فروش}
\label{\detokenize{user_guide/index:id7}}
\sphinxAtStartPar
در بخش مدیریت فروش، می‌توانید فروش‌ها را مدیریت کنید:
\begin{itemize}
\item {} 
\sphinxAtStartPar
\sphinxstylestrong{ثبت فروش جدید}: با کلیک بر روی دکمه "ثبت فروش جدید"، می‌توانید فروش جدید ثبت کنید

\item {} 
\sphinxAtStartPar
\sphinxstylestrong{جستجوی فروش}: با استفاده از فیلترهای مختلف، می‌توانید فروش‌ها را جستجو کنید

\item {} 
\sphinxAtStartPar
\sphinxstylestrong{ویرایش فروش}: با کلیک بر روی دکمه "ویرایش" در کنار هر فروش، می‌توانید اطلاعات فروش را ویرایش کنید

\item {} 
\sphinxAtStartPar
\sphinxstylestrong{حذف فروش}: با کلیک بر روی دکمه "حذف" در کنار هر فروش، می‌توانید فروش را حذف کنید

\item {} 
\sphinxAtStartPar
\sphinxstylestrong{مشاهده جزئیات فروش}: با کلیک بر روی شماره فروش، می‌توانید جزئیات فروش را مشاهده کنید

\end{itemize}


\subsection{مدیریت محصولات}
\label{\detokenize{user_guide/index:id8}}
\sphinxAtStartPar
در بخش مدیریت محصولات، می‌توانید محصولات را مدیریت کنید:
\begin{itemize}
\item {} 
\sphinxAtStartPar
\sphinxstylestrong{افزودن محصول جدید}: با کلیک بر روی دکمه "افزودن محصول جدید"، می‌توانید محصول جدید ایجاد کنید

\item {} 
\sphinxAtStartPar
\sphinxstylestrong{جستجوی محصول}: با استفاده از فیلترهای مختلف، می‌توانید محصولات را جستجو کنید

\item {} 
\sphinxAtStartPar
\sphinxstylestrong{ویرایش محصول}: با کلیک بر روی دکمه "ویرایش" در کنار هر محصول، می‌توانید اطلاعات محصول را ویرایش کنید

\item {} 
\sphinxAtStartPar
\sphinxstylestrong{حذف محصول}: با کلیک بر روی دکمه "حذف" در کنار هر محصول، می‌توانید محصول را حذف کنید

\item {} 
\sphinxAtStartPar
\sphinxstylestrong{مشاهده جزئیات محصول}: با کلیک بر روی نام محصول، می‌توانید جزئیات محصول را مشاهده کنید

\end{itemize}


\chapter{شاخص‌ها و جداول}
\label{\detokenize{index:id4}}\begin{itemize}
\item {} 
\sphinxAtStartPar
\DUrole{xref}{\DUrole{std}{\DUrole{std-ref}{genindex}}}

\item {} 
\sphinxAtStartPar
\DUrole{xref}{\DUrole{std}{\DUrole{std-ref}{modindex}}}

\item {} 
\sphinxAtStartPar
\DUrole{xref}{\DUrole{std}{\DUrole{std-ref}{search}}}

\end{itemize}



\renewcommand{\indexname}{فهرست}
\printindex
\end{document}